% sec_identifiability.tex — identifiability and black-box degradation
\section{Identifiability and Black-Box Degradation}\label{sec:identifiability}

Measurement claims are only as good as their identifiability analysis. This
section fixes what is knowable about autonomy, from what data, and with which
degradation rules when internal access is denied. Proofs are in
Appendix~\ref{app:identifiability}.

\subsection{White-box identification}

\begin{proposition}[White-box point identification, tiered by modal
commitment]\label{prop:whitebox}
Given white-box access---the full causal log of $S$, $O$, $A$, $I$, $Z$
unfolded to the horizon, the local laws of $\Gcal$, and intervention access
(H-0)---every channel estimand \eqref{eq:aorg}--\eqref{eq:aref} is
identified, but the premises differ by tier, and we state them apart rather
than compress them into ``functional of the law.''
\emph{Tier 1 (law functionals)}: $\Aorg$, $\Aepi$, $\Avia$ are functionals of
the joint law (directed informations, accuracy against the logged ground
truth, supply ratios).
\emph{Tier 2 (protocol-constructed modal objects)}: $\Aexec$ quantifies over
the accessible policy class $\Pi_t$ (a counterfactual object,
operationalized by calibration probes) and $\Aref$'s ensemble is constructed
by the measurement protocol; under those constructions the values are unique,
and without them the quantities are not defined, rather than unidentified.
\emph{Tier 3 (counterfactual)}: $\Lam$ \eqref{eq:lambda} is a path-specific
\emph{counterfactual}, not an interventional functional of the joint law:
computing it requires the structural model---the logged local laws and
exogenous noise terms that support the abduction step of the counterfactual
evaluation---so its identification premise is mechanism-level (white-box in
the strong sense), strictly stronger than joint-law access with interventions
(which pins interventional distributions but, absent structural assumptions,
does not pin cross-world counterfactuals). The MI proxy's certified domain is
narrower still (the AX3 conditions \textup{(V)}$\wedge$\textup{(A)}).
Estimation details and bias sources:
plug-in directed
information carries $O(\cdot/N)$ bias (simulation calibration is the
registered remedy); mutual-information estimation at scale uses
Kraskov-type $k$-NN or neural estimators
\citep{kraskov2004,belghazi2018}; empowerment is computable by variational
bounds \citep{mohamed2015} and Blahut--Arimoto-type iteration
\citep{cover_thomas2006}; offline change-point estimation covers the
registration layer \citep{truong2020,salge2013}.
\end{proposition}

Two per-channel notes complement the tier statement.
$\Aexec$ is exact only under the probe protocol that certifies the accessible
class (Section~\ref{sec:blackbox}).
$\Aref$'s measurement ensemble is \emph{constructed} by the protocol
(self-prediction tasks under the perturbation family): invariance of the
reading under admissible protocol re-designs is a reporting obligation, not
an automaticity. And the MI ratio for $\Lam$
presumes variation of the surgery history \emph{within the
ensemble}: in a single deployment with a constant surgery history the
empirical MI degenerates to $0$, and the randomized-audit design is the
proxy's only certified domain
(Section~\ref{sec:datareq}).

The measurement protocol is four-phase: a \emph{baseline period} records all
channel parameters and logs without intervention; a \emph{calibration period}
applies known interventions of known type and strength (under ethical review)
and fits the five channel response functions---this is also where unit
conversions (H2), the weights $w_i$ (by one of the three non-circular
determination protocols of Section~\ref{sec:representation}), the drift
tolerance $\eta$, and the strength ratio $\rho$ are estimated; the \emph{measurement
period} probes the components and composites $A$ with the elasticity
robustness band ($\varsigma$); a \emph{longitudinal design} repeats measurement to estimate
$dA/dt$ and the C1/C2 dynamical parameters.

\subsection{Per-channel data requirements}\label{sec:datareq}

\begin{center}\small
\begin{tabularx}{\textwidth}{@{}l >{\raggedright\arraybackslash}X@{}}
\toprule
Component & Data required \\
\midrule
$\Avia$ & resource supply/substitution inventory, maintenance-capacity audit \\
$\Aorg$ & full causal log $(S,O,A,I)$ \\
$\Agoal$ & behavior traces + endorsement protocol records \\
$\Aexec$ & environment interaction interface \\
$\Aepi$ & belief readout $+$ ground-truth outcomes (logged $E_t$;
verification probes for humans) \\
$\Aref$ & self-prediction tasks (ensemble, $H(\theta_\pi)>0$) + change
approval logs \\
$\Lam$ & full causal log (endorsement events + surgery history
$I^{\mathrm{surg},T}$; white-box) \\
$\AHI$ & parallel self-report (calibration track only) \\
\bottomrule
\end{tabularx}
\end{center}

\begin{center}\small
\resizebox{\textwidth}{!}{%
\begin{tabular}{@{}lll@{}}
\toprule
Component & Identification strategy/status & Main bias source \\
\midrule
$\Avia$ & direct audit (supply ratios; envelope capacities) & hidden
dependencies underestimated \\
$\Aorg$ & causally conditioned DI estimation
(closed loop needs intervention data) & closed-loop DI bias \\
$\Agoal$ & MaxEnt-IRL for $\Ueff$
(temperature fixed by protocol); $\Ustar$ elicited by structured deliberation
& IRL underdetermination (report intervals) \\
$\Aexec$ & variational empowerment over
the accessible policy class + capacity/utilization split & variance in
high-dimensional observations \\
$\Aepi$ & accuracy scoring
$\Pr[B_t=\kappa(E_t)]$ against the truth & belief
read-out error; misregistered coarse-graining $\kappa$ \\
$\Aref$ & self-probes + endorsement-consistency audit & probe
overfitting \\
$\Lam$ & law-level counterfactual \eqref{eq:lambda} (path-specific
$\mathrm{do}$ on the causal log); MI proxy certified only under
\textup{(V)}$\wedge$\textup{(A)} & black-box
unidentifiable in principle $\Rightarrow$ report $[\Lam_{\min},1]$ \\
$\AHI$ & $\Aphen-\Aobj$ & --- (warning signal only, never evidence) \\
\bottomrule
\end{tabular}}
\end{center}

\noindent The $\Ustar$ elicitation protocol (human substrate) is minimally
specified as \emph{structured deliberation}: a recorded sequence of pairwise
comparisons over the canonical parameter box $\Theta$ with the anchored
normalization $U(x_0)=0$, a fixed protocol metric $d$, and an auditable
decision record. Absent this record, $\Agoal$ degrades to the flagged
interval of \eqref{eq:agoal} rather than to a point value.

\subsection{Black-box degradation rules}\label{sec:blackbox}

Without internal access, only behavior--intervention responses are available.
What survives is stated with its certified scope: the certified direction of
each rule is a theorem under an explicit premise; the uncertified directions
are reported as heuristics and labeled as such.

\paragraph{$\Aorg$: a conditionally certified one-sided bound with a
reference lower end.}
By the randomized-probe lemma
(Lemma~\ref{lem:probe}, Appendix~\ref{app:prelim}), an exogenous randomized
probe makes $F_{ext}^{obs}=I(I^T;A^T)$ observable without access to $S$
(the lemma identifies the probe MI with the \emph{unconditioned} directed
information; the bypass premise of the proposition below is what promotes
the comparison to the causally conditioned estimand). The rule, with its two
premises stated \emph{in} the claim rather than around it:
\begin{proposition}[Conditionally certified upper bound for
$\Aorg$]\label{prop:certified-bound}
Assume \textup{(i)} bypass probe semantics, $I^t\ind S^t\given A^{t-1}$
(the probe is exogenous, non-anticipative, and acts outside the
belief-mediation path); and \textup{(ii)} internal-flow drift bounded by the
calibration-estimable tolerance, $\Fint\le\Fint^{base}(1+\eta)$,
$\eta\in[0,1]$. Then, with $F_{ext}^{obs}=I(I^T;A^T)$ the observable probe
flow,
\begin{equation}\label{eq:aorg-interval}
\Aorg\;\le\;
\min\Bigl(1,\frac{\Fint^{base}(1+\eta)}{\Fint^{base}(1+\eta)+F_{ext}^{obs}}\Bigr).
\end{equation}
Both premises are black-box unverifiable in general; the bound is therefore
\emph{conditionally} certified, and the BB-2 arm of
Section~\ref{sec:experiments} calibrates exactly what a premise violation
costs.
\end{proposition}
\begin{proof}[Proof sketch]
Under (i), $I(I^t;S^t\given A^{t-1})=0$, so the per-term identity displayed
in the scope paragraph below is non-negative termwise and
$F_{ext}^{obs}\le\Fext$; $\Aorg$ is decreasing in $\Fext$; with (ii),
substituting the drift bound gives \eqref{eq:aorg-interval}.
Remark~\ref{rem:probe-bridge} gives the bridge in full.
\end{proof}
\noindent The upper end of \eqref{eq:aorg-interval} is reported together with
the \emph{reference lower end}
$\Fint^{base}(1-\eta)/\bigl(\Fint^{base}(1-\eta)+F_{ext}^{obs}\bigr)$,
whose status is categorically different: it treats $F_{ext}^{obs}$ as exact,
and since $F_{ext}^{obs}$ may \emph{under}state the true flow, the truth can
fall below the lower end---it is a heuristic plug-in reference, not a
certified bound, and \eqref{eq:aorg-interval} must never be cited as an
interval-containment result. (Insufficient $\eta$ serves as the negative
control on premise (ii).) The band's
width's dependence on baseline length and $\eta$ is the information content
assessment; certified-direction results are reported in
Section~\ref{sec:experiments} (zero violations of the upper end in $3/3$
bypass arms, with widths reported).

\paragraph{Scope of the bound (the two premises, stated as obligations).}
(i) \emph{Probe semantics, and the direction of the bias}: the observable
$I(I^T;A^T)$ and the target
$\Fext=\sum_t I(I^t;A_t\given A^{t-1},S^t)$ differ in general. The per-term
identity
\[
I(I^t;A_t\given A^{t-1},S^t)-I(I^t;A_t\given A^{t-1})
= I(I^t;S^t\given A^{t-1},A_t)-I(I^t;S^t\given A^{t-1})
\]
shows that no direction holds in general. Two probe semantics exhaust the
audit practice. \emph{Bypass-semantics probes}---exogenous, non-anticipative,
and orthogonal to the agent's state given the action history,
$I^t\ind S^t\given A^{t-1}$ (a probe that acts outside the belief-mediation
path, as in the $\beta_\pi$-style script arms of BB-1): the second term of the
identity vanishes and every per-term difference is non-negative, so
$F_{ext}^{obs}\le\Fext$---the raw probe MI is then a \emph{lower} estimate of
the causally conditioned external flow, and the upper end of
\eqref{eq:aorg-interval} is certified. \emph{Mediation probes}---content
absorbed through $B\in S$, i.e.\ the
$\alpha$ class---reverse the direction: with $S=I\oplus N$, $A=S$ one has
$\Fext=0$ while $I(I;A)=0.531$~bit, so the plug-in \emph{over}estimates the
external flow and the certificate simply does not apply---its bypass premise
fails on this family (the truth lies above the plug-in band). The certified
scope of \eqref{eq:aorg-interval} is therefore the bypass family---not
``non-adversarial influence'' in general, since honest $\alpha$ influence is
mediated by construction. The BB-1 anchors (zero upper-end violations in
$3/3$ bypass arms) exercise
bypass arms and do not cover the mediation case. (ii) \emph{Adversarial
adaptation}: the bypass premise requires $I$ exogenous, non-anticipative, and
state-orthogonal; an adaptive adversary whose injection timing correlates
with the agent's state violates it, and the naive exogenous plug-in
misestimates the flow. The benchmark's BB-2 arm calibrates exactly this
failure---the certified bound's premise is violated in $3/3$ adaptive arms,
the naive exogenous plug-in then overestimates the external flow by $+0.17$
to $+0.90$ nats, and the bound no longer holds there; exogenous control arms
retain it. Auditing
adaptive adversaries requires causal-layer access---reinforcing
Corollary~\ref{cor:logs-necessary} from the estimation side.

\paragraph{$\Aexec$: one-sided bound with protocol dependence.} Under a
\emph{passive observation} protocol (policy-driven actions), the observable
proxy overestimates: $\Aexec^{obs}\ge\Aexec$ (state entropy mixes in). Under an
\emph{independent probe} protocol ($A\ind S$), the direction reverses:
conditioning denoises, $I(A;O)\le I(A;O\given S)$. The direction of the bound
is therefore protocol-relative, and reporting the protocol label is a
disclosure obligation---an empirical nuance our benchmark confirms in both
directions (Section~\ref{sec:experiments}).

\paragraph{$\Lam$: impossibility as a theorem.}

\begin{proposition}[Implanted endorsement is black-box indistinguishable]
\label{prop:unidentifiable}
There exist two models $\mathcal{M}_{endo}$ (endogenous endorsement converging
to $\Ustar$) and $\mathcal{M}_{impl}$ (a $\beta$ surgery at $t=0$ installing
$G_0=\Ustar=U_{adv}$ with endorsement thenceforth the identity) such that
$P(O^T,A^T)$ is pointwise identical (and the $\Aphen$ laws agree), while
$\Lam(\mathcal{M}_{endo})=1$ and $\Lam(\mathcal{M}_{impl})=0$. Hence every
black-box functional assigns them the same $A$: historicality is in principle
not black-box identifiable. We state the epistemic anatomy of the claim
exactly: the \emph{black-box half} (pointwise identity of the observable
laws) is proved by construction and needs no convention; the white-box value
$\Lam(\mathcal{M}_{impl})=0$ is the registered certification convention of
AX3 (a degenerate endorsement process under a non-empty surgery history
certifies nothing)---under the counterfactual definition \eqref{eq:lambda}
alone the value would depend on the construction's counterfactual target,
which the construction leaves unspecified. The convention is stated in AX3
and invoked here; with it, the separation is exact.
\end{proposition}
\begin{proof}[Proof sketch]
Construction: take one behavioral norm (policy $\pi^*$); $\mathcal{M}_{endo}$
converges to $\Ustar$ through endorsement; $\mathcal{M}_{impl}$ has $G_0$
installed at $t=0$ with $\Endorse\equiv\mathrm{id}$. Both generate the same
$(O,A)$ marginal law from the same $\pi^*$ and $o$; they differ only at the
unobserved node $I^\beta$. Appendix~\ref{app:identifiability}.
\end{proof}

\begin{corollary}[Full causal logging is an architectural
necessity]\label{cor:logs-necessary}
(i) Systematic manipulation cases with a procedurally perfect counterfeit
history \citep{mele2008}
(as discussed by \citealp{fischer_ravizza1998}) and deceptive alignment
\citep{hubinger2019}
are two instances of one mathematical fact: a perfect procedural counterfeit
is behaviorally equivalent to the genuine article. (ii) Black-box mode reports
$\Lam\in[\Lam_{\min},1]$ with $\Lam_{\min}$ available only from causal
logs/supply-chain provenance audit (cf.\ \citealp{capability_permission2026}).
(iii) The full-causal-logging requirement of our benchmark platform
(Section~\ref{sec:experiments}) is therefore an identifiability
\emph{necessity}, not an engineering convenience.
\end{corollary}

A mediation-analysis alternative (natural direct effects
\citep{pearl2001}: $\Fext=0\iff\mathrm{NDE}(I\to A)=0$) requires sequential
ignorability \citep{vanderweele2015} and cross-world
counterfactuals---assumptions strictly stronger than the certified-bound
route; we take \eqref{eq:aorg-interval} as primary and NDE as a sensitivity
direction (numerics pending; Section~\ref{sec:discussion}).
